\chapter{Truecolour, palette, and degradation}
\label{ch:colour}

\begin{aside}
A modern terminal claims to support 16 million colours. The
terminal you SSHed into last week claims sixteen. \maya{}'s job is
to make your app look as good as possible on each, without writing
the app twice.
\end{aside}

\section{Three colour modes}

Terminal colour history is a stack of additions, each one keeping
the previous:

\begin{enumerate}[leftmargin=*]
  \item \textbf{Sixteen.} The original ANSI palette: eight base,
    eight bright. Emitted as \code{ESC [ 30..37 m} and \code{ESC
    [ 90..97 m} for foreground; \code{40..47} and \code{100..107}
    for background.
  \item \textbf{256.} An xterm extension: 16 base + 216 RGB cube
    (6x6x6) + 24 greys. Emitted as \code{ESC [ 38 ; 5 ; N m}.
  \item \textbf{Truecolour.} 24-bit RGB. \code{ESC [ 38 ; 2 ; R ;
    G ; B m}. Standardised by ISO 8613-6 in 1994; widely
    supported since about 2015.
\end{enumerate}

\maya{} represents every colour as 24-bit RGB internally and
degrades at emit time based on the detected terminal capabilities.

\section{Detection}

The \code{COLORTERM} environment variable is the de-facto switch:

\begin{itemize}[leftmargin=*]
  \item \code{COLORTERM=truecolor} or \code{=24bit}: truecolour.
  \item Otherwise check \code{TERM}: \code{xterm-256color},
    \code{screen-256color}, etc. imply 256.
  \item Otherwise check \code{TERM}: \code{xterm}, \code{vt100},
    etc. imply 16.
  \item \code{TERM=dumb} or unset means no colour at all.
\end{itemize}

\begin{lstlisting}
ColorMode detect_colour_mode() {
    if (auto* ct = std::getenv("COLORTERM"))
        if (std::string_view(ct) == "truecolor"
         || std::string_view(ct) == "24bit")
            return ColorMode::Truecolour;
    auto* term = std::getenv("TERM");
    if (!term) return ColorMode::None;
    std::string_view t = term;
    if (t.contains("256color")) return ColorMode::Indexed256;
    if (t == "dumb")            return ColorMode::None;
    return ColorMode::Indexed16;
}
\end{lstlisting}

\section{Degradation}

When the requested colour exceeds the terminal's capability,
\maya{} maps it down. There are two strategies:

\begin{itemize}[leftmargin=*]
  \item \textbf{Nearest.} For each input RGB, walk the palette
    and pick the entry with smallest Euclidean distance in RGB
    space.
  \item \textbf{Perceptual.} Convert to CIELAB and pick by
    $\Delta E$. This handles eye sensitivity better but costs a
    matrix multiply per pixel.
\end{itemize}

\maya{} uses nearest, with the 16-colour palette using a
hand-tuned mapping that splits the cube differently from a naive
nearest-neighbour (the result looks better for typical UI
accents).

\section{The 16-colour palette}

The base 16 are not fixed RGB values. They are slots, and the
user's terminal theme picks the actual colour. The standard names
(\code{red}, \code{bright\_red}, etc.) map to slots, not pixels.

This is a feature: a user who chose a Solarized theme expects
your ``red'' to be Solarized's red. \maya{} does not override; it
emits the slot and lets the terminal do its job.

\begin{warning}
Hardcoded RGB ``corrections'' to the base 16 break this. Resist
the urge. If you want a specific shade, use truecolour and
accept the degradation on 16-colour terminals.
\end{warning}

\section{The 256-colour cube}

The 6x6x6 RGB cube uses non-linear levels:
\code{\{0, 95, 135, 175, 215, 255\}}. The gap between 0 and 95 is
huge; the gaps after are 40 each. The non-linearity matches human
brightness perception roughly but not well.

A truecolour-to-cube map:

\begin{lstlisting}
int channel_to_cube(uint8_t v) {
    static constexpr int levels[6] =
        {0, 95, 135, 175, 215, 255};
    int best = 0, best_d = 1000;
    for (int i = 0; i < 6; ++i) {
        int d = std::abs(int(v) - levels[i]);
        if (d < best_d) { best_d = d; best = i; }
    }
    return best;
}
\end{lstlisting}

Then the palette index is \code{16 + 36*r + 6*g + b}.

The grey ramp (indices 232--255) is a separate 24-step linear
grey. If R = G = B, mapping to the grey ramp gives a more
faithful tone than the cube, which has only 6 grey samples on
its diagonal.

\section{Default colours}

\code{ESC [ 39 m} resets foreground to default; \code{49} resets
background. \maya{} emits these when a cell explicitly uses
\code{Color::default}. Saving the default-reset means the
terminal background shows through, which is what users on
transparent or themed terminals expect.

\section{NO\_COLOR}

\begin{funfact}
Jason Boice proposed the \code{NO\_COLOR} environment variable
in 2017. Set it to any value and conforming programs must emit
no ANSI colour codes. Over 250 programs honour it today,
including \code{ls}, \code{grep}, \code{git}, and \maya{}.
\end{funfact}

\maya{} checks \code{NO\_COLOR} at startup and clamps to
\code{ColorMode::None} if set. Style attributes (bold, italic,
underline) still emit; only colour is suppressed.

\section{Dim, bold, and SGR interaction}

The \code{ESC [ 1 m} (bold) and \code{ESC [ 2 m} (dim) codes
have an unfortunate interaction with the base 16 palette: on
many terminals, bold makes the bright variant fire even if the
slot is the dim one. So \code{red + bold} looks identical to
\code{bright\_red}.

\maya{} does \emph{not} compensate. The behaviour is
user-configured and varies; trying to ``fix'' it would surprise
users who tuned their terminal.

\section{Colour spaces and gamma}

Terminal RGB is sRGB, gamma 2.2 approximately. Computing
gradients in sRGB looks wrong (the middle of a gradient is too
dark). For accurate gradients, convert to linear, interpolate,
and convert back:

\begin{lstlisting}
float to_linear(uint8_t v) {
    float f = v / 255.0f;
    return f <= 0.04045f
         ? f / 12.92f
         : std::pow((f + 0.055f) / 1.055f, 2.4f);
}
\end{lstlisting}

This matters for the gradient widget and for any element that
blends colours. \maya{} provides \code{lerp\_srgb} as a helper
that does the round trip.

\section{Hex helpers}

Most users think in hex. \maya{} accepts hex literals as
compile-time strings:

\begin{lstlisting}
auto label = text("hello") | fg<"#FF6432">;
\end{lstlisting}

The parser is \code{constexpr}, so a typo (\code{"FF6"} or
\code{"\#ZZZZZZ"}) is a compile error, not a runtime one.

\section{Theme tokens}

For consistency across a UI, name your colours:

\begin{lstlisting}
namespace theme {
  constexpr auto accent     = "\#FF6432";
  constexpr auto background = "\#1E1E1E";
  constexpr auto muted      = "\#8E8E8E";
}
\end{lstlisting}

Then \code{fg<theme::accent>} reads the role, not the hex.
Changing the theme is one edit in one place.

\section{Recap}

\begin{recap}
\begin{itemize}[leftmargin=*]
  \item Three modes: 16, 256, truecolour.
  \item Detect via \code{COLORTERM} and \code{TERM}.
  \item \maya{} stores RGB internally and degrades on emit.
  \item Honour \code{NO\_COLOR}.
  \item Base 16 are slots, not pixels; let the user's theme
    pick.
  \item Hex literals are validated at compile time.
\end{itemize}
\end{recap}

\section*{Workshop}

\begin{enumerate}[leftmargin=*]
  \item Print the result of \code{detect\_colour\_mode()} at
    startup. Run with \code{TERM=xterm}, \code{TERM=xterm-256color},
    \code{COLORTERM=truecolor}, and \code{NO\_COLOR=1}.
  \item Build a gradient widget that fades from one hex to
    another. Compare sRGB-linear and naive lerp side by side.
  \item Add a \code{--theme} flag to your app that swaps the
    \code{theme::} namespace at startup. Persist the choice in
    a config file.
\end{enumerate}
